\chapter{A Difference Appears}

\section{The first fact}

The instrument's first act is not a measurement. Before it can weigh, count, or read anything, it does something
smaller and stranger: it tells a \emph{this} from a \emph{that}. That capacity --- to \emph{distinguish} --- is the
machine's ground floor, and the construction names it exactly that: the distinguishable.

The test is austere. The machine carries a single running index --- a \emph{level} --- and two things are
distinguishable precisely when they sit at different levels. Nothing else is consulted. This is the identity of
indiscernibles that Leibniz~\cite{leibniz1686} stated three centuries ago, run backwards: if you cannot tell two
things apart, they are one; the instant you can, a difference exists. The level-index itself is the ordinary
hierarchy of \emph{universes} from type theory~\cite{martin1970}, pressed into one humble service --- to give the
word ``different'' a mechanical meaning, so that a machine and not a person can decide it.

Notice what has \emph{not} happened. Nothing is measured. There is no quantity, no unit, no dial --- and, the
construction is emphatic about this, no storage: the difference allocates nothing. It is prior to counting. The
machine must be able to say \emph{these two are not the same} before it can say how many, or how much, and that
saying is free.

This is where any instrument has to begin. You cannot read a scale whose marks you cannot resolve into distinct
things; resolution is logically upstream of measurement --- the same point Nyquist~\cite{nyquist1928} and
Shannon~\cite{shannon1948} would later make about signals, made here about symbols. The whole book is the story of
what the machine builds on this one act, and, at the very end, of the single thing it builds \emph{toward} and still
cannot pin. It starts here: a difference, told apart, not yet counted.

\section{The symbol as a label}

Where does the machine put a difference once it has one? Nowhere --- and that is the point. The distinction is not
written into a cell to be looked up later; it \emph{is} the label. When the machine tells two things apart, what it
produces is exactly the running level-index at which they part ways. The difference and its record are one object.

This is why it can be said, without cheating, that nothing has been allocated. A conventional program that wanted to
remember ``these two differ'' would reserve a bit, a byte, a slot --- storage whose cost you pay and whose
bookkeeping you keep. The machine here does not. It uses the one index it was already carrying --- the level it must
track anyway to keep its own descriptions coherent --- and reads the difference straight off it. The label is free
because the ledger was already open.

There is an old idea underneath this, older than any computer. A symbol, as Peirce~\cite{peirce1878} and then the
whole of information theory would insist, means nothing in isolation; it means only \emph{by contrast with the
symbols it is not}. Shannon~\cite{shannon1949} made it quantitative --- a symbol carries information exactly to the
degree it could have been another --- but the qualitative core is Saussure's~\cite{saussure1916}: value is
difference. The machine takes that literally. Its symbols \emph{are} their differences; the label is the
distinction, not a name pinned on afterward.

So the first thing the instrument hands you is not a measurement but a \emph{position}: which level, relative to
what. Everything it will later call a number, a charge, a mass, a cost, is built by iterating this one move --- take
a difference, let it be its own label --- and paying, each time, only the turn.
